\documentclass[10pt,a4paper]{amsart}
\usepackage{amsmath,amssymb,amsthm,graphicx}
% Retained minimal preamble for the Stacks Project chapter more-algebra.tex (More on Algebra).
% The chapter file pinned-original-source.tex carries only \input{preamble} on its first line, so no
% preamble text lives inside the chapter itself.  The Stacks Project uses one single shared preamble for
% every chapter; the snapshot stages it identically as sources/stacks-chow/preamble.tex and
% sources/stacks-etale-cohomology/preamble.tex.  The block below is copied mechanically, line for line,
% from sources/stacks-chow/preamble.tex lines 14-18 (the Xy-pic setup) and 154-217 (the native theorem
% declarations together with the mathematical macros).  Nothing was added, removed or redefined; the
% wrapper class and the amsmath/amssymb/amsthm/graphicx base packages are supplied by the assembler.
% Chapter: more-algebra.tex -- More on Algebra.
% For commutative diagrams we use Xy-pic
\usepackage[all]{xy}

% We use 2cell for 2-commutative diagrams.
\xyoption{2cell}

% Theorem environments.
%
\theoremstyle{plain}
\newtheorem{theorem}[subsection]{Theorem}
\newtheorem{proposition}[subsection]{Proposition}
\newtheorem{lemma}[subsection]{Lemma}

\theoremstyle{definition}
\newtheorem{definition}[subsection]{Definition}
\newtheorem{example}[subsection]{Example}
\newtheorem{exercise}[subsection]{Exercise}
\newtheorem{situation}[subsection]{Situation}

\theoremstyle{remark}
\newtheorem{remark}[subsection]{Remark}
\newtheorem{remarks}[subsection]{Remarks}

\numberwithin{equation}{subsection}

% Macros
%
\def\lim{\mathop{\mathrm{lim}}\nolimits}
\def\colim{\mathop{\mathrm{colim}}\nolimits}
\def\Spec{\mathop{\mathrm{Spec}}}
\def\Hom{\mathop{\mathrm{Hom}}\nolimits}
\def\Ext{\mathop{\mathrm{Ext}}\nolimits}
\def\SheafHom{\mathop{\mathcal{H}\!\mathit{om}}\nolimits}
\def\SheafExt{\mathop{\mathcal{E}\!\mathit{xt}}\nolimits}
\def\Sch{\mathit{Sch}}
\def\Mor{\mathop{\mathrm{Mor}}\nolimits}
\def\Ob{\mathop{\mathrm{Ob}}\nolimits}
\def\Sh{\mathop{\mathit{Sh}}\nolimits}
\def\NL{\mathop{N\!L}\nolimits}
\def\CH{\mathop{\mathrm{CH}}\nolimits}
\def\proetale{{pro\text{-}\acute{e}tale}}
\def\etale{{\acute{e}tale}}
\def\QCoh{\mathit{QCoh}}
\def\Ker{\mathop{\mathrm{Ker}}}
\def\Im{\mathop{\mathrm{Im}}}
\def\Coker{\mathop{\mathrm{Coker}}}
\def\Coim{\mathop{\mathrm{Coim}}}

% Boxtimes
%
\DeclareMathSymbol{\boxtimes}{\mathbin}{AMSa}{"02}

%
% Macros for moduli stacks/spaces
%
\def\QCohstack{\mathcal{QC}\!\mathit{oh}}
\def\Cohstack{\mathcal{C}\!\mathit{oh}}
\def\Spacesstack{\mathcal{S}\!\mathit{paces}}
\def\Quotfunctor{\mathrm{Quot}}
\def\Hilbfunctor{\mathrm{Hilb}}
\def\Curvesstack{\mathcal{C}\!\mathit{urves}}
\def\Polarizedstack{\mathcal{P}\!\mathit{olarized}}
\def\Complexesstack{\mathcal{C}\!\mathit{omplexes}}
% \Pic is the operator that assigns to X its picard group, usage \Pic(X)
% \Picardstack_{X/B} denotes the Picard stack of X over B
% \Picardfunctor_{X/B} denotes the Picard functor of X over B
\def\Pic{\mathop{\mathrm{Pic}}\nolimits}
\def\Picardstack{\mathcal{P}\!\mathit{ic}}
\def\Picardfunctor{\mathrm{Pic}}
\def\Deformationcategory{\mathcal{D}\!\mathit{ef}}
\usepackage{cleveref}
\crefname{originallemma}{original lemma}{original lemmas}
\Crefname{originallemma}{Original lemma}{Original lemmas}
\makeatletter
\newlabel{more-algebra-lemma-check-tor-dimension-modulo-nilpotent}{{0H75}{source}{}{sourcepointer-1}{}}
\newlabel{more-algebra-lemma-check-tor-dimension-modulo-nilpotent@cref}{{[originallemma][0][]0H75}{[1][0][]source}{}{}{}}
\makeatother
\crefname{originallemma}{original lemma}{original lemmas}
\Crefname{originallemma}{Original lemma}{Original lemmas}
\makeatletter
\newlabel{lemma-pull-tor-amplitude}{{066L}{source}{}{sourcepointer-2}{}}
\newlabel{lemma-pull-tor-amplitude@cref}{{[originallemma][0][]066L}{[1][0][]source}{}{}{}}
\makeatother
\title{Tor amplitude modulo a nilpotent ideal agrees with tor amplitude over the ring}
\author{The Stacks Project Authors}
\date{stacks-project:more-algebra.tex:local-snapshot-20261012; GNU FDL 1.2 or later}
\begin{document}
\maketitle
\setcounter{subsection}{1}
\noindent\textit{Editorial note.} All delivered mathematical text is the author's own, copied from the pinned Stacks Project chapter file more-algebra.tex; no statement or proof step was rewritten, reordered or removed.\par\medskip
\noindent\textit{Editorial note.} The chapter file pinned-original-source.tex carries only \textbackslash{}input\{preamble\} on its first line, so no preamble text lives inside the chapter. The retained file sources/stacks-more-algebra/retained-preamble.tex is a mechanical, line-for-line copy of the project's single shared preamble as staged in the snapshot (Xy-pic setup plus theorem declarations and mathematical macros). Nothing was added, removed or redefined, and the excerpt serves only as preamble, never as a span coordinate.\par\medskip
\noindent\textit{Editorial note.} Both counted spans start at the \textbackslash{}begin line of their environment, so the statement and the author's proof are delivered whole; surrounding prose of the section is deliberately outside the counted spans and is carried as standard prerequisites instead.\par\medskip
\noindent\textit{Editorial note.} The citation fields (source\_id, source\_version, authors, source\_url, license.name, source\_title) are identical to those of every other card this lane has produced from this chapter, so the batch selector treats them as one source.\par\medskip
\noindent\textit{Allowed standard prerequisites (editorial).}\begin{itemize}
\item Tor amplitude of a complex and its behaviour under base change along a ring map. the author cites Lemma lemma-pull-tor-amplitude (tag 066L) at proof line 17315. Gives the easy direction of the equivalence.
\item Derived tensor product of a complex with a module, computed by resolving the module. used at proof lines 17320-17335. Supplies the object whose cohomology is being bounded.
\item Distinguished triangle attached to a short exact sequence of modules and the resulting long exact sequence in cohomology. used at proof lines 17330-17344. Is the tool that transfers the amplitude bound from the two outer terms to the middle one.
\end{itemize}\medskip
\ifdefined\hypertarget\hypertarget{sourcepointer-1}{}\fi
\noindent\textit{Source reference.} Original lemma 0H75: focused-seeds/AGP/sources/stacks-snapshot-20261011/tags.txt, source lines 21445--21445. The lemma's own canonical tag, quoted so that readers can resolve the label inside the delivered text; it points at the counted statement itself and carries no proof content.\par\medskip
\ifdefined\hypertarget\hypertarget{sourcepointer-2}{}\fi
\noindent\textit{Source reference.} Original lemma 066L: focused-seeds/AGP/sources/stacks-snapshot-20261011/tags.txt, source lines 7940--7940. Navigation pointer to a statement the counted proof refers to but does not contain; nothing load-bearing is delegated to it.\par\medskip
\begin{lemma}
\label{lemma-check-tor-dimension-modulo-nilpotent}
Let $R' \to R$ be a surjective ring map whose kernel is a nilpotent ideal.
Let $K' \in D(R')$ and set $K = K' \otimes_{R'}^\mathbf{L} R$.
Let $a, b \in \mathbf{Z}$.
Then $K$ has tor amplitude in $[a, b]$ if and only if $K'$ does.
\end{lemma}

\begin{proof}
One direction follows from Lemma \ref{lemma-pull-tor-amplitude}.
For the other, assume $K$ has tor amplitude in $[a, b]$ and let
$M'$ be an $R'$-module. We have to show that $K' \otimes_{R'}^\mathbf{L} M'$
has nonzero cohomology only for degrees contained in the interval $[a, b]$.

\medskip\noindent
Let $I = \Ker(R' \to R)$. Then $I^n = 0$ for some $n$. If $IM' = 0$,
then we can view $M'$ as an $R$-module and argue as follows
$$
K' \otimes_{R'}^\mathbf{L} M' =
K' \otimes_{R'}^\mathbf{L} (R \otimes_R^\mathbf{L} M') =
(K' \otimes_{R'}^\mathbf{L} R) \otimes_R^\mathbf{L} M' =
K \otimes_R^\mathbf{L} M'
$$
which has nonvanishing cohomology only in the interval $[a, b]$
by assumption on $K$. If $I^{t + 1}M' = 0$, then we consider the
short exact sequence
$$
0 \to IM' \to M' \to M'/IM' \to 0
$$
By induction on $t$ we have that both $K' \otimes_{R'}^\mathbf{L} IM'$
and $K' \otimes_{R'}^\mathbf{L} M'/IM'$ have nonzero cohomology only for
degrees in the interval $[a, b]$. Then the distinguished triangle
$$
K' \otimes_{R'}^\mathbf{L} IM' \to
K' \otimes_{R'}^\mathbf{L} M' \to
K' \otimes_{R'}^\mathbf{L} M'/IM' \to
(K' \otimes_{R'}^\mathbf{L} IM')[1]
$$
proves the same is true for $K' \otimes_{R'}^\mathbf{L} M'$ as desired.
\end{proof}

\end{document}
